\documentclass[11pt]{article}
\usepackage[margin=1.08in]{geometry}
\usepackage{amsmath,amssymb,amsthm,mathtools,booktabs,array,microtype}
\usepackage[T1]{fontenc}
\usepackage{lmodern}
\usepackage[hidelinks]{hyperref}
\hypersetup{
  pdftitle={Unconditional low-multiplicity profiles for zeros of the Riemann zeta function},
  pdfauthor={Zach Waddle},
  pdfsubject={Unconditional bounds for bounded zero multiplicity from a multiplicity-aware rank--trace inequality},
  pdfkeywords={Riemann zeta function, zero multiplicity, pair correlation, rank--trace inequality}
}
\usepackage{xcolor}
\usepackage{enumitem}
\usepackage{needspace}
\usepackage{longtable}
\usepackage{fancyhdr}
\usepackage{lastpage}

\newtheorem{theorem}{Theorem}[section]
\newtheorem{proposition}[theorem]{Proposition}
\newtheorem{lemma}[theorem]{Lemma}
\newtheorem{corollary}[theorem]{Corollary}
\newtheorem{remark}[theorem]{Remark}
\newtheorem{definition}[theorem]{Definition}

\newcommand{\N}{\mathcal N}
\newcommand{\D}{\mathcal D}
\newcommand{\kMT}{\kappa_{\mathrm{MT}}}
\newcommand{\ind}{\mathbf 1}
\newcommand{\R}{\mathbb R}
\newcommand{\C}{\mathbb C}
\DeclareMathOperator{\tr}{tr}
\DeclareMathOperator{\rank}{rank}

\pagestyle{fancy}
\fancyhf{}
\lhead{Unconditional low-multiplicity profiles}
\rhead{Z. Waddle}
\cfoot{\thepage\ of \pageref{LastPage}}
\setlength{\headheight}{14pt}

\title{\textbf{Unconditional low-multiplicity profiles for zeros of the Riemann zeta function}\\[0.4em]
\large Sharp consequences of a multiplicity-aware rank--trace inequality}
\author{Zach Waddle}
\date{Private draft for expert review -- August 12, 2026 (version 0.2)}

\begin{document}
\maketitle

\begin{abstract}
A recent unconditional two-trace argument of Claude/Anthropic proves that at least
$0.6725007\ldots$ of the nontrivial zeros of the Riemann zeta function are simple and on the critical line and at least $0.8362503\ldots$ are distinct.  Its formal development contains a stronger one-parameter, multiplicity-aware rank--trace inequality valid for every real parameter $c>0$.  We optimize the complete scalar family rather than only its $c=2$ and $c=3$ specializations.

For each fixed integer $r\geq 1$, we give unconditional lower bounds for (i) the fraction of the zero multiset supported on locations of multiplicity at most $r$, and (ii) the fraction of distinct zero locations having multiplicity at most $r$.  For example, at least
\[
0.887620008173354\ldots
\]
of the zero multiset is supported on locations of multiplicity at most three, while at least
\[
0.969319059130202\ldots
\]
of the distinct zero locations have multiplicity at most three.  The corresponding bounds for multiplicity at most four are $0.906350006811128\ldots$ and $0.979753104026191\ldots$.

We also construct extremal abstract multiplicity distributions showing that every displayed bound is sharp if one retains only the total zero count and the full scalar rank--trace family.  Thus any further improvement from the same analytic two-trace input must use geometric information discarded by those scalar inequalities.  The new work is finite-dimensional algebra and integer multiplicity bookkeeping; no new explicit-formula or prime-correlation estimate is required.
\end{abstract}

\paragraph{Status and provenance.}
The sole analytic input is the unconditional two-trace theorem and multiplicity-aware Lean development of \cite{Claude2026}.  The present version was audited directly against commit \path{3635e74826a4c1fcece7d1cd2b6fa75e43a00510}: the arbitrary-$c$ matrix inequality is \path{rank_trace_mult_k_le}, and the required perturbation estimate is the coefficient-generic theorem \path{ctr_sub_frobSq_perturb}.  The scalar optimization and all displayed constants were independently regenerated by exact rational interval arithmetic and symbolic checks.  The new corollaries have not yet been added to the Lean tree or externally refereed.  This is therefore a private research draft, not a claim of completed peer review.

\section{Introduction}

Let $\rho=\beta+i\gamma$ run over the distinct nontrivial zeros of $\zeta(s)$, and let $m_\rho\geq1$ denote multiplicity.  Write
\[
N(T_1,T_2)=\sum_{T_1<\gamma\leq T_2}m_\rho,
\qquad
N_d(T_1,T_2)=\#\{\rho:T_1<\gamma\leq T_2\}.
\]
The recent argument \cite{Claude2026} gives, with an optimized Montgomery--Taylor window,
\[
\liminf_{T\to\infty}\frac{N_0^s(T,2T)}{N(T,2T)}
\geq 2-\kMT
=0.672500703679411\ldots,
\]
and
\[
\liminf_{T\to\infty}\frac{N_d(T,2T)}{N(T,2T)}
\geq \frac{3-\kMT}{2}
=0.836250351839705\ldots,
\]
where
\begin{equation}\label{eq:kappa}
\kMT
=\frac{1}{2}+\frac{1}{\sqrt2}\cot\!\left(\frac{1}{\sqrt2}\right)
=1.327499296320588354265620209\ldots.
\end{equation}
Here $N_0^s$ counts simple critical-line zeros.

The higher-level multiplicity problem has previously been studied through $n$-level correlations.  Assuming the generalized Riemann hypothesis, Gon\c{c}alves, de Laat, and Leijenhorst \cite{GDL2025} prove that at least $0.9614$ of the zeta-zero multiset has multiplicity at most two and at least $0.9787$ has multiplicity at most three.  Explicit unconditional bounds for zeros of a prescribed large multiplicity are given by Simoni\v{c}, Trudgian, and Turnage-Butterbaugh \cite{STTB2022}.  The present note addresses a different regime: fixed low multiplicities, without RH, using only the two-trace input of \cite{Claude2026}.

A targeted search of the accessible primary literature did not locate an earlier unconditional fixed-$r$ hierarchy with the formulas below, especially the distinct-location profile and the sharpness statement for the complete scalar family.  This is a provisional priority assessment rather than a claim of first publication; a specialist bibliography check remains part of the review process.

For $r\geq1$, define the \emph{low-multiplicity multiset count}
\begin{equation}\label{eq:Mrdef}
M_{\leq r}(T_1,T_2)
:=\sum_{\substack{\rho\ \mathrm{distinct}\\T_1<\gamma\leq T_2\\m_\rho\leq r}}m_\rho
\end{equation}
and the \emph{low-multiplicity distinct count}
\begin{equation}\label{eq:Drdef}
D_{\leq r}(T_1,T_2)
:=\#\{\rho:T_1<\gamma\leq T_2,\ m_\rho\leq r\}.
\end{equation}
Thus $M_{\leq1}=D_{\leq1}$ is the number of simple zeros at all horizontal locations, while $M_{\leq r}$ counts every copy of a zero whose location has multiplicity at most $r$.

Our constants use
\begin{equation}\label{eq:Bdef}
B:=(3-2\sqrt2)(\kMT-1)
=0.056189995913322830066962070\ldots.
\end{equation}
The exact enclosure in Section~\ref{sec:verification} gives
\begin{equation}\label{eq:basicbounds}
1<\kMT<\frac43,
\qquad
0<B<\frac1{15}.
\end{equation}
In particular, $r-1-rB>0$ for every $r\geq3$, so all denominators below are positive.

\begin{theorem}[Multiset low-multiplicity profile]\label{thm:multiset}
For the nontrivial zeros of $\zeta(s)$,
\begin{align}
\liminf_{T\to\infty}\frac{M_{\leq1}(T,2T)}{N(T,2T)}
&\geq 2-\kMT,\label{eq:M1}\\
\liminf_{T\to\infty}\frac{M_{\leq2}(T,2T)}{N(T,2T)}
&\geq \frac{3-\kMT}{2},\label{eq:M2}
\end{align}
and, for every fixed integer $r\geq3$,
\begin{equation}\label{eq:Mr}
\liminf_{T\to\infty}\frac{M_{\leq r}(T,2T)}{N(T,2T)}
\geq 1-B\frac{r+1}{r-1}.
\end{equation}
The same inequalities hold for cumulative counts on $(0,T]$ and for every fixed primitive Dirichlet $L$-function.
\end{theorem}

\begin{theorem}[Distinct-location multiplicity profile]\label{thm:distinct}
For the nontrivial zeros of $\zeta(s)$,
\begin{align}
\liminf_{T\to\infty}\frac{D_{\leq1}(T,2T)}{N_d(T,2T)}
&\geq 1-\frac{\kMT-1}{3-\kMT},\label{eq:D1}\\
\liminf_{T\to\infty}\frac{D_{\leq2}(T,2T)}{N_d(T,2T)}
&\geq 1-\frac{\kMT-1}{2(4-\kMT)},\label{eq:D2}
\end{align}
and, for every fixed integer $r\geq3$,
\begin{equation}\label{eq:Dr}
\liminf_{T\to\infty}\frac{D_{\leq r}(T,2T)}{N_d(T,2T)}
\geq 1-\frac{B}{r-1-rB}.
\end{equation}
The same inequalities hold cumulatively and for every fixed primitive Dirichlet $L$-function.
\end{theorem}

The first line of Theorem~\ref{thm:distinct} is also an elementary consequence of the distinct-zero lower bound in \cite{Claude2026}.  The higher lines use the full one-parameter rank--trace family and are the main new content.

\begin{table}[ht]
\centering
\caption{Rigorous lower bounds, rounded downward.  The middle column counts copies in the zero multiset; the final column counts distinct locations.}
\label{tab:values}
\begin{tabular}{@{}rcc@{}}
\toprule
$r$ & $M_{\leq r}/N$ & $D_{\leq r}/N_d$\\
\midrule
1 & 0.672500703679411 & 0.804185854391506\\
2 & 0.836250351839705 & 0.938727930759812\\
3 & 0.887620008173354 & 0.969319059130202\\
4 & 0.906350006811128 & 0.979753104026191\\
5 & 0.915715006130015 & 0.984891304068350\\
6 & 0.921334005721348 & 0.987949456853123\\
8 & 0.927755719540013 & 0.991422003329087\\
10& 0.931323338328160 & 0.993340918495729\\
20& 0.937895267674748 & 0.996856714757252\\
\bottomrule
\end{tabular}
\end{table}

\section{The master scalar inequality}

We isolate the exact consequence of the analytic input that is used in the remainder of the paper.  For $c>0$ and $m\geq0$, put
\begin{equation}\label{eq:kc}
k_c(m):=c^2-\bigl((c-m)_+\bigr)^2
=\begin{cases}
2cm-m^2,&m\leq c,\\
c^2,&m\geq c.
\end{cases}
\end{equation}
For $m\geq0$ one has
\begin{equation}\label{eq:kcbounds}
0\leq k_c(m)\leq c^2.
\end{equation}
Indeed, $0\leq(c-m)_+\leq c$.

For the actual interval $(T,2T]$, let $x_m(T)$ be the number of critical-line zero locations of multiplicity $m$, and let $y_m(T)$ be the number of off-line reflected pairs whose two members each have multiplicity $m$.  Thus
\begin{equation}\label{eq:countidentity}
N(T,2T)=\sum_{m\geq1}m x_m(T)+2\sum_{m\geq1}m y_m(T).
\end{equation}

For $0<\lambda<1$, let
\[
c_\lambda^*=
\frac{\sqrt2\tan(\lambda/\sqrt2)}
{1+(\lambda/\sqrt2)\tan(\lambda/\sqrt2)},
\qquad
\kappa_\lambda=(c_\lambda^*)^{-1}.
\]
The Montgomery--Taylor optimization in \cite{Claude2026} gives $\kappa_\lambda\to\kMT$ as $\lambda\to1^-$.  The following epsilon-form is the precise interface needed below.

\begin{proposition}[Arbitrary-$c$ master inequality]\label{prop:master}
For every fixed $c>0$ and every $\varepsilon>0$, there exists $T_0=T_0(c,\varepsilon)$ such that, for all $T\geq T_0$,
\begin{equation}\label{eq:master}
(2c-\kMT-\varepsilon)N(T,2T)
\leq
\sum_{m\geq1}k_c(m)x_m(T)+c^2\sum_{m\geq1}y_m(T).
\end{equation}
\end{proposition}

\begin{proof}
Fix $c>0$ and $\varepsilon>0$.  Choose $\lambda\in(0,1)$ sufficiently close to one that
\[
\kappa_\lambda<\kMT+\frac{\varepsilon}{4}.
\]
Use the Montgomery--Taylor window at this fixed bandwidth.  Let $\widehat A$ be the normalized compression formed from the zeros in the enlarged height window, and write $\widehat G=\widehat A+\widehat E$.

The zero-block normalization and Poisson lemmas in
\path{Zeta23/ZeroSide/Mult.lean} identify the on-line part with the matrix model used by the
kernel-checked theorem \path{rank_trace_mult_k_le} in
\path{Zeta23/ZeroSide/RankTraceMult.lean} and verify its vector-norm and positive-index hypotheses.  Applying that theorem gives, for every $c>0$,
\begin{equation}\label{eq:ranktracegeneric}
2c\,\tr\widehat A-\|\widehat A\|_F^2
\leq
\sum_{\rho\ \mathrm{on\ line\ in}\ I'}k_c(m_\rho)+c^2p(I'),
\end{equation}
where $p(I')$ is the number of reflected off-line pairs in the enlarged window.  No independence of the evaluation vectors is assumed.

The kernel-checked coefficient-generic perturbation theorem \path{ctr_sub_frobSq_perturb} in \path{Zeta23/Assembly/SeamMult.lean} gives
\begin{equation}\label{eq:perturb}
2c\,\tr\widehat G-\|\widehat G\|_F^2
-B_T\bigl(2c+2\|\widehat G\|_F+B_T\bigr)
\leq
2c\,\tr\widehat A-\|\widehat A\|_F^2,
\end{equation}
where $B_T\to0$ for fixed $\lambda$.  The trace theorem for the same window gives
\[
\tr\widehat G=N(T,2T)+o_\lambda(N(T,2T)),
\qquad
\|\widehat G\|_F^2\leq
\bigl(\kappa_\lambda+o_\lambda(1)\bigr)N(T,2T).
\]
Consequently $\|\widehat G\|_F=O_\lambda(N^{1/2})$, and the perturbation term in \eqref{eq:perturb} is $o_{c,\lambda}(N)$.

The two boundary strips have total zero multiplicity $O(T^{1/2}\log T)=o(N)$.  By \eqref{eq:kcbounds}, every boundary on-line location costs at most $c^2$, while every boundary pair also costs $c^2$; their combined scalar charge is therefore at most $c^2$ times the boundary multiplicity and is $o_c(N)$.  Removing the boundary terms from \eqref{eq:ranktracegeneric}, and then taking $T$ sufficiently large, yields
\[
\left(2c-\kappa_\lambda-\frac{3\varepsilon}{4}\right)N(T,2T)
\leq
\sum_{m\geq1}k_c(m)x_m(T)+c^2\sum_{m\geq1}y_m(T).
\]
The choice of $\lambda$ implies \eqref{eq:master}.
\end{proof}

\begin{remark}[What is and is not reused]\label{rem:reuse}
The source tree currently specializes the final Montgomery--Taylor multiplicity endgame to $c=2$ and $c=3$.  Proposition~\ref{prop:master} does not assume an unproved new analytic estimate: it combines two already kernel-checked arbitrary-coefficient lemmas with the already formalized fixed-bandwidth trace and tail asymptotics, then performs the elementary boundary-charge estimate above.  A future Lean extension must still package this combination as one generic theorem, but no new explicit formula or prime-correlation input is required.
\end{remark}

The rest of the paper is a sharp finite-dimensional consequence of Proposition~\ref{prop:master}.

\section{The multiset hierarchy}

We first record a general majorization device.

\begin{lemma}[Multiset majorant]\label{lem:multmaj}
Fix $r\geq1$ and suppose $c,\alpha,\beta>0$ satisfy, for every integer $m\geq1$,
\begin{align}
k_c(m)&\leq \alpha m+\beta m\ind_{m\leq r},\label{eq:onmaj}\\
c^2&\leq 2\alpha m+2\beta m\ind_{m\leq r}.\label{eq:pairmaj}
\end{align}
Then
\begin{equation}\label{eq:multgeneric}
\liminf_{T\to\infty}\frac{M_{\leq r}(T,2T)}{N(T,2T)}
\geq \frac{2c-\kMT-\alpha}{\beta}.
\end{equation}
\end{lemma}

\begin{proof}
Sum \eqref{eq:onmaj} over critical-line points and \eqref{eq:pairmaj} over reflected pairs.  By \eqref{eq:countidentity}, the $\alpha$-terms sum exactly to $\alpha N$ and the $\beta$-terms to $\beta M_{\leq r}$.  Proposition~\ref{prop:master} gives, for every $\varepsilon>0$ and all sufficiently large $T$,
\[
(2c-\kMT-\varepsilon)N\leq \alpha N+\beta M_{\leq r}.
\]
Divide by $N$, take the lower limit, and then let $\varepsilon\downarrow0$.
\end{proof}

\begin{proof}[Proof of Theorem~\ref{thm:multiset}]
For $r=1$, take
\[
c=2,\qquad \alpha=2,\qquad \beta=1.
\]
Then $k_2(1)=3=2+1$, while $k_2(m)=4\leq2m$ for $m\geq2$; the pair inequality is immediate.  Lemma~\ref{lem:multmaj} gives $2-\kMT$.

For $r=2$, take
\[
c=3,\qquad \alpha=3,\qquad \beta=2.
\]
Here $k_3(1)=5$, $k_3(2)=8$, and $k_3(m)=9$ for $m\geq3$, so the two pointwise inequalities again hold.  This gives $(3-\kMT)/2$.

Now let $r\geq3$ and set
\begin{equation}\label{eq:Cdef}
C:=2+\sqrt2,
\qquad
\alpha_r:=\frac{C^2}{r+1},
\qquad
\beta_r:=\frac{C^2(r-1)}{2(r+1)}.
\end{equation}
The balancing identity
\begin{equation}\label{eq:balance}
C^2=2(2C-1)
\end{equation}
implies
\[
\alpha_r+\beta_r=2C-1=\frac{C^2}{2}.
\]
For $m\leq r$ and $m\leq C$, equation \eqref{eq:kc} gives
\[
k_C(m)=m(2C-m)\leq m(2C-1)=m(\alpha_r+\beta_r).
\]
For $m\leq r$ and $m\geq C$, one has
\[
k_C(m)=C^2\leq \frac{C^2}{2}m=m(\alpha_r+\beta_r).
\]
For $m\geq r+1$, since $r+1\geq4>C$,
\[
k_C(m)=C^2=\alpha_r(r+1)\leq\alpha_rm.
\]
This proves the on-line majorant.  For a reflected pair, if $m\leq r$ then
\[
2m(\alpha_r+\beta_r)=mC^2\geq C^2,
\]
and if $m\geq r+1$ then $2\alpha_rm\geq2C^2>C^2$.  Hence Lemma~\ref{lem:multmaj} applies and yields
\[
\liminf_{T\to\infty}\frac{M_{\leq r}(T,2T)}{N(T,2T)}
\geq
\frac{2C-\kMT-\alpha_r}{\beta_r}.
\]
Direct simplification, using $(2+\sqrt2)^{-2}=3/2-\sqrt2$, gives
\[
\frac{2C-\kMT-\alpha_r}{\beta_r}
=1-(3-2\sqrt2)(\kMT-1)\frac{r+1}{r-1},
\]
which is \eqref{eq:Mr}.

The cumulative result follows by dyadic summation exactly as in \cite{Claude2026}.  The fixed primitive Dirichlet $L$-function case uses the corresponding trace theorem from that paper; the zero-side scalar argument is unchanged.
\end{proof}

\section{The profile among distinct locations}

Put
\[
H_r(T_1,T_2):=N_d(T_1,T_2)-D_{\leq r}(T_1,T_2),
\]
the number of distinct locations of multiplicity at least $r+1$.

\begin{lemma}[Distinct-location majorant]\label{lem:distmaj}
Fix $r\geq1$.  Suppose $c,\alpha,\beta,\gamma>0$ satisfy $\alpha=2c-\kMT$ and, for every integer $m\geq1$,
\begin{align}
k_c(m)&\leq \alpha m+\beta-\gamma\ind_{m\geq r+1},\label{eq:diston}\\
c^2&\leq 2\alpha m+2\beta-2\gamma\ind_{m\geq r+1}.\label{eq:distpair}
\end{align}
Then
\begin{equation}\label{eq:distgeneric}
\limsup_{T\to\infty}\frac{H_r(T,2T)}{N_d(T,2T)}\leq\frac{\beta}{\gamma}.
\end{equation}
\end{lemma}

\begin{proof}
Summing the pointwise inequalities and using \eqref{eq:countidentity} gives
\[
\sum_m k_c(m)x_m+c^2\sum_m y_m
\leq \alpha N+\beta N_d-\gamma H_r.
\]
For every $\varepsilon>0$, Proposition~\ref{prop:master} and $\alpha=2c-\kMT$ imply, for all sufficiently large $T$,
\[
\gamma H_r\leq\beta N_d+\varepsilon N.
\]
The distinct-zero theorem in \cite{Claude2026} gives $N_d\geq\delta N$ eventually for some fixed $\delta>0$.  Hence
\[
\limsup_{T\to\infty}\frac{H_r}{N_d}
\leq\frac{\beta}{\gamma}+\frac{\varepsilon}{\gamma\delta}.
\]
Let $\varepsilon\downarrow0$.
\end{proof}

\begin{proof}[Proof of Theorem~\ref{thm:distinct}]
For $r=1$, take
\[
c=2,
\quad \alpha=4-\kMT,
\quad \beta=\kMT-1,
\quad \gamma=3-\kMT.
\]
For $m=1$ the on-line inequality is equality.  For $m\geq2$, the right side is increasing in $m$ and equals $k_2(2)=4$ at $m=2$.  The pair inequality has nonnegative slack.  Lemma~\ref{lem:distmaj} gives \eqref{eq:D1}.

For $r=2$, take
\[
c=3,
\quad \alpha=6-\kMT,
\quad \beta=\kMT-1,
\quad \gamma=2(4-\kMT).
\]
Again the on-line inequality is equality at $m=1$ and $m=3$, and has positive slack at $m=2$; the pair inequality is weaker.  This gives \eqref{eq:D2}.

For $r\geq3$, use $C=2+\sqrt2$ and set
\begin{equation}\label{eq:distparams}
\alpha:=2C-\kMT,
\qquad
\beta:=\kMT-1,
\qquad
\gamma:=\alpha(r+1)+\beta-C^2.
\end{equation}
For $m=1$, equation \eqref{eq:balance} gives
\[
k_C(1)=2C-1=\alpha+\beta.
\]
For $m=2,3$ one checks exactly that
\begin{equation}\label{eq:smallslack}
\alpha m+\beta-k_C(m)=(m-1)\bigl(m-(\kMT-1)\bigr)\geq0.
\end{equation}
For $4\leq m\leq r$, $k_C(m)=C^2$.  Since $\alpha>0$, it is enough to check $m=4$; here
\[
4\alpha+\beta-C^2=4C-3\kMT+1>0,
\]
using $\kMT<4/3$.  For $m\geq r+1$, the right side of \eqref{eq:diston} is increasing in $m$ and equals $C^2$ at $m=r+1$ by the definition of $\gamma$.  Thus \eqref{eq:diston} holds.  The pair inequality is equality for a simple pair and has positive slack otherwise.

Lemma~\ref{lem:distmaj} therefore gives
$\limsup_{T\to\infty}H_r(T,2T)/N_d(T,2T)\leq\beta/\gamma$.  Writing $B=(3-2\sqrt2)(\kMT-1)$ and simplifying,
\[
\frac{\beta}{\gamma}=\frac{B}{r-1-rB},
\]
which proves \eqref{eq:Dr}.  Cumulative and Dirichlet variants follow as before.
\end{proof}

\section{Sharpness of the complete scalar family}

We now show that the constants above cannot be improved using only Proposition~\ref{prop:master} for all $c>0$ and the total-count identity.  This is a sharpness statement for the scalar relaxation, not a claim that the extremal distributions occur among actual zeta zeros.

Normalize total multiplicity to one.  Let $x_m\geq0$ denote the density of on-line locations of multiplicity $m$, and $y_m\geq0$ the density of off-line pairs of multiplicity $m$.  The abstract feasibility conditions are
\begin{equation}\label{eq:feasible}
\sum_{m\geq1}m x_m+2\sum_{m\geq1}m y_m=1
\end{equation}
and
\begin{equation}\label{eq:feasiblec}
2c-\kMT\leq\sum_{m\geq1}k_c(m)x_m+c^2\sum_{m\geq1}y_m
\qquad(c>0).
\end{equation}

\begin{theorem}[Sharpness]\label{thm:sharp}
For each fixed $r\geq1$, the lower bounds in Theorems~\ref{thm:multiset} and \ref{thm:distinct} are optimal over all nonnegative distributions satisfying \eqref{eq:feasible}--\eqref{eq:feasiblec}.
\end{theorem}

\begin{proof}
For $r=1$, take
\[
x_1=2-\kMT,
\qquad
x_2=\frac{\kMT-1}{2},
\qquad
x_m=y_m=0\ \text{otherwise}.
\]
Then \eqref{eq:feasible} holds.  If $0<c\leq1$, the slack in \eqref{eq:feasiblec} is decreasing and is nonnegative at $c=1$; for $1\leq c\leq2$ it equals
\[
\frac{\kMT-1}{2}(c-2)^2;
\]
for $c\geq2$ it is identically zero.  The low multiset mass is $2-\kMT$, and the low fraction among distinct locations is
\[
\frac{x_1}{x_1+x_2}=1-\frac{\kMT-1}{3-\kMT}.
\]

For $r=2$, take
\[
x_1=\frac{3-\kMT}{2},
\qquad
x_3=\frac{\kMT-1}{6},
\qquad
x_m=y_m=0\ \text{otherwise}.
\]
For $1\leq c\leq3$, the slack is
\[
\frac{\kMT-1}{6}(c-3)^2,
\]
while the ranges $c\leq1$ and $c\geq3$ are handled as above.  This realizes both $r=2$ bounds.

Now suppose $r\geq3$.  Put
\[
d:=\kMT-1,
\qquad
A:=3-2\sqrt2,
\]
and define
\begin{equation}\label{eq:sharpabc}
a:=1-d(\sqrt2-1),
\qquad
h:=\frac{Ad}{r-1},
\qquad
b:=\frac{1-a-(r+1)h}{2}.
\end{equation}
\Needspace{8\baselineskip}
All three numbers are nonnegative.  Indeed, $a>0$ because $0<d<1$ and $0<\sqrt2-1<1$, while $h>0$ and
\[
2b=d\left((\sqrt2-1)-(3-2\sqrt2)\frac{r+1}{r-1}\right)
\geq d(5\sqrt2-7)>0,
\]
since $(r+1)/(r-1)\leq2$ and $\sqrt2>7/5$.  Take
\[
x_1=a,
\qquad
x_{r+1}=h,
\qquad
y_1=b,
\]
with all other variables zero.  By construction,
\[
a+(r+1)h+2b=1.
\]
Let
\[
S(c):=a k_c(1)+h k_c(r+1)+bc^2-(2c-\kMT).
\]
For $1\leq c\leq r+1$, direct expansion using \eqref{eq:sharpabc} gives the exact square
\begin{equation}\label{eq:square}
S(c)=(h+b)(c-(2+\sqrt2))^2\geq0.
\end{equation}
For $0<c\leq1$,
\[
S(c)=(a+h+b)c^2-2c+\kMT.
\]
Since $a+h+b\leq1$, this is decreasing on $(0,1]$ and is nonnegative at $c=1$ by \eqref{eq:square}.  For $c\geq r+1$, differentiation and the total-count identity give
\[
S'(c)=2b(c-2)\geq0,
\]
so this range is minimized at $c=r+1$, where \eqref{eq:square} again applies.  Hence the distribution satisfies every constraint \eqref{eq:feasiblec}.

Its low multiset mass is
\[
a+2b
=1-Ad\frac{r+1}{r-1},
\]
which is exactly \eqref{eq:Mr}.  Its distinct total and high-multiplicity distinct mass are
\[
D=a+h+2b=1-\frac{Adr}{r-1},
\qquad
H_r=h=\frac{Ad}{r-1},
\]
so
\[
\frac{H_r}{D}=\frac{Ad}{r-1-rAd},
\]
which is equality in \eqref{eq:Dr}.  This proves sharpness.
\end{proof}

\begin{remark}[Why $2+\sqrt2$ appears]
For $r\geq3$, a simple on-line point has scalar charge $k_c(1)=2c-1$ when $c\geq1$, while each point in a simple off-line pair carries half of the pair charge, $c^2/2$.  Balancing those two worst cases gives
\[
2c-1=\frac{c^2}{2},
\]
whose relevant solution is $c=2+\sqrt2$.  The irrational parameter is therefore structural rather than a numerical accident.
\end{remark}

\section{Consequences and comparison}

Theorems~\ref{thm:multiset} and \ref{thm:distinct} imply, for example,
\[
\limsup_{T\to\infty}
\frac{N(T,2T)-M_{\leq3}(T,2T)}{N(T,2T)}
\leq0.112379991826645\ldots,
\]
so at most $11.2380\%$ of the zero multiset may be supported on locations of multiplicity at least four.  Among distinct locations, at most $3.0681\%$ may have multiplicity at least four.

For $j\geq2$, define the exact-multiplicity multiset and distinct counts
\[
E_j(T_1,T_2):=\sum_{\substack{T_1<\gamma\leq T_2\\m_\rho=j}}m_\rho,
\qquad
E_{d,j}(T_1,T_2):=\#\{\rho:T_1<\gamma\leq T_2,\ m_\rho=j\}.
\]
Taking $r=j-1$ in the two profile theorems gives the following exact-multiplicity upper bounds.

\begin{corollary}[Exact multiplicity profile]\label{cor:exact}
One has
\begin{align*}
\limsup_{T\to\infty}\frac{E_2(T,2T)}{N(T,2T)}
&\leq \kMT-1,\\
\limsup_{T\to\infty}\frac{E_3(T,2T)}{N(T,2T)}
&\leq \frac{\kMT-1}{2},\\
\limsup_{T\to\infty}\frac{E_j(T,2T)}{N(T,2T)}
&\leq B\frac{j}{j-2}\qquad(j\geq4),
\end{align*}
and
\begin{align*}
\limsup_{T\to\infty}\frac{E_{d,2}(T,2T)}{N_d(T,2T)}
&\leq \frac{\kMT-1}{3-\kMT},\\
\limsup_{T\to\infty}\frac{E_{d,3}(T,2T)}{N_d(T,2T)}
&\leq \frac{\kMT-1}{2(4-\kMT)},\\
\limsup_{T\to\infty}\frac{E_{d,j}(T,2T)}{N_d(T,2T)}
&\leq \frac{B}{j-2-(j-1)B}\qquad(j\geq4).
\end{align*}
Every bound is sharp within the scalar information model of Theorem~\ref{thm:sharp}.
\end{corollary}

\begin{proof}
The exact-multiplicity-$j$ set is contained in the set of locations of multiplicity at least $j$, so the inequalities follow from Theorems~\ref{thm:multiset} and \ref{thm:distinct} with $r=j-1$.  The sharpness distributions for $r=j-1$ place all of their high-multiplicity mass at multiplicity exactly $j$.
\end{proof}

For fixed $r$, these are unconditional two-trace bounds.  They are weaker numerically than the GRH-dependent higher-correlation results of \cite{GDL2025}; for instance, that work obtains $0.9614$ at multiplicity at most two and $0.9787$ at multiplicity at most three.  The distinction is that no RH or higher-level correlation input is used here.

Corollary~\ref{cor:exact} also complements the unconditional exact-multiplicity theorem of \cite{STTB2022}.  Their estimate is uniform in $j$ and decays exponentially, becoming effective relative to elementary simple-zero bounds for extremely large multiplicity.  The present estimates are fixed-$j$ consequences of two traces and are aimed at the low-multiplicity regime; in particular, the multiset upper bound $Bj/(j-2)$ is not asserted uniformly when $j$ grows with $T$.

As $r\to\infty$, the multiset lower bound tends to
\[
1-B=0.943810004086677\ldots,
\]
whereas the distinct-location lower bound tends to one with tail
\[
1-\frac{D_{\leq r}}{N_d}
\leq \frac{B}{r}+O\!\left(\frac1{r^2}\right).
\]
Thus the scalar two-trace information forces high multiplicities to occupy a vanishing fraction of distinct zero locations, although it still permits a fixed fraction of the zero multiset to escape to multiplicities growing with height.

\paragraph{Separate derivative problem.}
The analogous optimization for zeros of $\xi'$ is not used in this paper.  It requires profile-specific analytic integration beyond the scalar bookkeeping above and is therefore kept in a separate research package rather than folded into the present theorem claim.

\section{Exact verification, source audit, and formalization path}\label{sec:verification}

The accompanying standard-library Python certificate uses exact rational arithmetic.  The identity
\[
\kMT=\frac12+\frac{\cos(1/\sqrt2)}{\sin(1/\sqrt2)/(1/\sqrt2)}
\]
allows both numerator and denominator to be enclosed by alternating Taylor series with rational terms:
\[
\cos(1/\sqrt2)=\sum_{n\geq0}\frac{(-1)^n}{2^n(2n)!},
\qquad
\frac{\sin(1/\sqrt2)}{1/\sqrt2}
=\sum_{n\geq0}\frac{(-1)^n}{2^n(2n+1)!}.
\]
The certificate encloses $\kMT$, $B$, every table entry, and every exact-multiplicity upper bound by rational intervals.  All printed lower bounds are rounded downward; all printed upper bounds are rounded upward.  A second script uses exact symbolic algebra to verify the parameter simplifications, the square identities in the sharpness proof, and the pointwise majorants over their complete integer ranges.

The source audit at commit \path{3635e74826a4c1fcece7d1cd2b6fa75e43a00510} records the following dependency chain:
\Needspace{12\baselineskip}
\begin{description}[leftmargin=1.3em,style=nextline,font=\normalfont\ttfamily]
\item[ZeroSide.Mult block normalization and Poisson lemmas]
identify the normalized zero compression and verify the hypotheses of the generic matrix theorem;
\item[ZeroSide.RankTraceMult.rank\_trace\_mult\_k\_le]
arbitrary-$c$ inequality \eqref{eq:ranktracegeneric};
\item[Assembly.ctr\_sub\_frobSq\_perturb]
arbitrary trace coefficient in \eqref{eq:perturb};
\item[ThmD.tracesBoundsD\_concrete and tendsto\_cRatio\_concrete]
fixed-$\lambda$ Montgomery--Taylor trace asymptotics;
\item[eventually\_tailPackageD and local-count inputs]
$B_T\to0$ and $O(T^{1/2}\log T)$ boundary multiplicity.
\end{description}

A complete Lean extension now has a finite scope:
\begin{enumerate}[leftmargin=2em]
\item define $M_{\leq r}$, $D_{\leq r}$, and the exact-multiplicity counts;
\item package Proposition~\ref{prop:master} by combining the existing generic rank--trace and perturbation declarations with the existing Theorem-D trace/tail endgame;
\item formalize the elementary scalar majorants and sharpness identities;
\item instantiate at $c=2$, $c=3$, and $c=2+\sqrt2$, then reuse the existing dyadic and fixed-Dirichlet-$L$ summation layers;
\item run the pinned complete build and the repository's \texttt{\#print axioms} audit.
\end{enumerate}
No new explicit-formula, prime-side, or zero-localization theorem is needed.  Lean and Lake are not installed in the present runtime, so the extension supplied with this draft is a source-level blueprint rather than a compiled formal proof.

\section{Limitations}

The results constrain zero multiplicities but do not prove the Riemann hypothesis. Sharpness in Theorem~\ref{thm:sharp} is sharpness of a scalar relaxation: the extremal distributions need not arise from the analytic evaluation vectors of the zeta function. Further improvement from the same two-trace input would require additional geometric or arithmetic information.

\section*{Acknowledgments and disclosure}
The analytic input and its formal development are due to the author of \cite{Claude2026} and the teams acknowledged there. Zach Waddle directed the present project. OpenAI's GPT-5.6 Pro assisted with derivation, checking, source inspection, exact-certificate generation, and manuscript preparation. Independent specialist review and a compiled Lean extension remain outstanding.

\begin{thebibliography}{99}
\small
\bibitem{Claude2026} Claude, \emph{More than two thirds of the zeros of the Riemann zeta function lie on the critical line}, preprint, 2026; accompanying Lean repository \texttt{anthropics/zeta-23-lean}, commit \path{3635e74826a4c1fcece7d1cd2b6fa75e43a00510}.
\bibitem{GDL2025} F. Gon\c{c}alves, D. de Laat, and N. Leijenhorst, \emph{Multiplicity of nontrivial zeros of primitive $L$-functions via higher-level correlations}, Math. Comp. 94 (2025), 2041--2058.
\bibitem{STTB2022} A. Simoni\v{c}, T. Trudgian, and C. L. Turnage-Butterbaugh, \emph{Some explicit and unconditional results on gaps between zeroes of the Riemann zeta-function}, Trans. Amer. Math. Soc. 375 (2022), 3239--3265.
\bibitem{Montgomery1973} H. L. Montgomery, \emph{The pair correlation of zeros of the zeta function}, Proc. Sympos. Pure Math. 24 (1973), 181--193.
\bibitem{CGG1998} J. B. Conrey, A. Ghosh, and S. M. Gonek, \emph{Simple zeros of the Riemann zeta-function}, Proc. London Math. Soc. 76 (1998), 497--522.
\end{thebibliography}


\end{document}
